% !TEX root = ../main.tex
\section{System Invariants}
\label{sec:invariants}

Sections~\ref{sec:architecture} and~\ref{sec:canonical-form} argue from
architecture. This section argues from contract, which is the form a reviewer can
check without agreeing with me about mechanism.

If you were specifying a component whose assertions you intended to rely on, you
would write down invariants. Two groups are needed, and the distinction between
them is the paper's practical claim in compressed form. The first group states
properties of assertion \emph{behaviour}: things you would assume without thinking
about them from a database, a rule engine, or a calculation service. All of them
fail, and none can be made to hold by work on $M$. The second states properties the
\emph{system} must supply. None of them holds by default, and all of them can be
engineered.

\subsection{Behavioural invariants: what does not hold}

\begin{table}[ht]
  \centering
  \small
  \begin{tabular}{@{}>{\raggedright\arraybackslash}p{0.5cm}>{\raggedright\arraybackslash}p{2.9cm}>{\raggedright\arraybackslash}p{4.1cm}>{\raggedright\arraybackslash}p{1.1cm}>{\raggedright\arraybackslash}p{4.4cm}@{}}
    \toprule
    & Invariant & Statement & Holds & What breaks without it \\
    \midrule
    I1 & \textbf{Form invariance}
       & $x \sim_c x'$ implies equivalent decision-relevant outcome
       & No
       & Tests do not generalise past the wording tested \\
    \addlinespace
    I2 & \textbf{Closure}
       & $\Out_S(p)$ and $\Out_S(p \rightarrow q)$ imply $\Out_S(q)$
       & No
       & Conclusions the system's own premises entail are unavailable \\
    \addlinespace
    I3 & \textbf{Decomposition}
       & $\Out_S(p \wedge q)$ implies $\Out_S(p)$
       & No
       & Summaries assert what their parts deny \\
    \addlinespace
    I4 & \textbf{Consistency}
       & Not both $\Out_S(p)$ and $\Out_S(\neg p)$
       & No
       & Two users, or one user twice, get opposite answers \\
    \addlinespace
    I5 & \textbf{Self-report fidelity}
       & Claims about what it knows track what it asserts
       & No
       & Confidence cannot be used to route or gate \\
    \addlinespace
    I6 & \textbf{Correctness}
       & $\Out_S(p)$ implies $p$
       & No
       & Output cannot be relied on unchecked \\
    \addlinespace
    I7 & \textbf{Repeatability}
       & Same input, same session, same assertion
       & No
       & A verified answer is not a verified behaviour \\
    \bottomrule
  \end{tabular}
  \caption{Behavioural invariants. I6 is the one everybody already knows fails, and
    it is the least interesting: nobody deploys these systems believing output is
    true by construction. I1 and I7 are the ones that invalidate validation.}
  \label{tab:invariants-behaviour}
\end{table}

One reading to head off first. These are invariants, so \emph{holds} is a universal
claim and a single counterexample defeats it. ``No'' in the third column therefore
means the property is not guaranteed, not that behaviour is always inconsistent. In
practice these systems satisfy I1--I4 on most inputs most of the time, which is
precisely why the failures are hard to catch in testing and why a specification cannot
rely on them. An invariant that holds usually is not an invariant; it is a tendency,
and nothing downstream can be built on it.

Three further observations, since the pattern matters more than the entries.

\paragraph{I6 failing is priced in; I1 failing is not.} Every deployment I have
seen assumes the system is sometimes wrong and builds review around it. Almost none
assumes that a reviewed, approved, tested behaviour will not reproduce under a
rephrasing. The first is a known unknown that review handles. The second
invalidates the review.

\paragraph{I2--I4 share one cause.} These are not three independent defects. There
is no component whose function is to maintain the set of asserted claims---nothing
that closes it under consequence, decomposes conjunctions, or checks it for
contradiction. By Section~\ref{sec:what-computed} the asserted ``set'' is not a
maintained data structure at all; it is whatever successive forward passes emit.
Absent such a component, I2--I4 have nothing to be true of.

\paragraph{I5 failing has a specific and costly consequence.} A great deal of
proposed architecture routes on model-reported confidence: escalate when the model
says it is unsure, proceed when it says it is confident. I5's failure means the
routing signal is produced by the same process as the content and is not a
measurement of it. Verbalised confidence can be trained to better aggregate
calibration \citep{lin2022}, and models have some capacity to evaluate their own
outputs \citep{kadavath2022}; neither delivers the per-claim reliability a gate
requires. Aggregate calibration and per-claim boundary recognition are different
properties, and only the second supports a decision about \emph{this} answer.

That last point generalises, and it is the shape of the whole argument: a property
that holds on average over a distribution is not a property you can gate on for an
individual case. Accuracy is such a property, as
Proposition~\ref{prop:independence} shows. So is calibration.

\subsection{System invariants: what must be supplied}

The constructive claim. These are conditions on $S$, not on $M$. None holds by
default; each is achievable by architecture; and together they are what must be
true for a response to cross from ordinary generative assistance into consequential
decision support.

\begin{table}[ht]
  \centering
  \small
  \begin{tabular}{@{}>{\raggedright\arraybackslash}p{0.7cm}>{\raggedright\arraybackslash}p{3.3cm}>{\raggedright\arraybackslash}p{4.6cm}>{\raggedright\arraybackslash}p{4.4cm}@{}}
    \toprule
    & Invariant & Requirement & Control that supplies it \\
    \midrule
    I8 & \textbf{Source identity}
       & Every cited source has stable identity and established authenticity
       & Controlled corpus, ingestion authentication \\
    \addlinespace
    I9 & \textbf{Temporal validity}
       & Source version and effective date are recorded and compared against $t$
       & Version metadata, corpus release pinning \\
    \addlinespace
    I10 & \textbf{Context binding}
       & $c$ is explicit before assertion, or explicitly marked unknown
       & Intake resolution, typed request objects \\
    \addlinespace
    I11 & \textbf{Applicability}
       & Evidence is filtered on scope, not similarity alone
       & Metadata-aware retrieval, scope predicates \\
    \addlinespace
    I12 & \textbf{Traceability}
       & Each consequential claim links to the evidence supporting it
       & Claim decomposition, claim-level citation \\
    \addlinespace
    I13 & \textbf{Support verification}
       & The cited evidence is checked to support the claim at the strength asserted
       & Entailment checks, policy logic, human review by tier \\
    \addlinespace
    I14 & \textbf{Reproducibility}
       & The assertion can be re-derived from recorded inputs and versions
       & Snapshot pinning, deterministic replay \\
    \addlinespace
    I15 & \textbf{Authorization}
       & Reliance is permitted at the applicable tier, with duties separated
       & Tiered policy, approval gates, access control \\
    \addlinespace
    I16 & \textbf{Reconstructability}
       & The full decision path is retained for the required period
       & Audit record, retention policy \\
    \addlinespace
    I17 & \textbf{Controlled failure}
       & Unmet conditions produce qualification, refusal, or escalation---not an
         unqualified answer
       & Output gating outside $M$ \\
    \addlinespace
    I18 & \textbf{Sufficiency}
       & The evidence set is checked for coverage of every condition the claim requires
       & Required-field or checklist verification against the rule \\
    \addlinespace
    I19 & \textbf{Defeater search}
       & A scoped search for contrary or undermining authority of equal or higher
         standing is performed and recorded
       & Governed defeater query; institution's precedence order \\
    \bottomrule
  \end{tabular}
  \caption{System invariants. I8--I9 are ordinary data governance. I10--I11 are the
    context work that Section~\ref{sec:cell-v} shows is most often absent.
    I12--I13 are what distinguishes grounding from citation. I14--I16 are what an
    audit actually asks for. I17 is what makes the others enforceable rather than
    advisory. I18--I19 correspond to the coverage and no-defeater conditions of
    Definition~\ref{def:grounding}, and exist so that every condition in the definition
    has a system obligation that establishes it.}
  \label{tab:invariants-system}
\end{table}

Two points about the second table.

\paragraph{I17 is the load-bearing one.} Without a controlled-failure path, every
other system invariant degrades into a preference. A system that detects missing
applicability and answers anyway has not enforced I11; it has measured it. And by
Section~\ref{sec:missing} the gate cannot live inside $M$, because generated
refusal is a continuation rather than a branch. I17 is therefore an architectural
requirement, not a prompt.

\paragraph{Proportionality, not maximalism.} Not every use requires all ten.
Section~\ref{sec:tiers} makes the controls proportional to the consequence of
reliance, and a drafting assistant needs none of I13--I19. What does not follow is
the converse: a system cannot claim auditability because it reports an accuracy
score, emits citations, or carries a generic disclaimer. Those are compatible with
failing every invariant in Table~\ref{tab:invariants-system}.
